\subsection[Eichpunkt: Schmelzpunkt von H2O]{Eichpunkt: Schmelzpunkt von \ce{H2O}}
Die Temperatur des Wasser-Eis-Gemisches kann auch unter 0\unit{\degreeCelsius} betragen. Eigentlich ist es ein System aus zwei Subsystemen Wasser und Eis mit unterschiedlicher Temperatur \(T_1>T_2\). Dieses will nun in einen stationären Zustand des thermodynamischen Gleichgewichts übergehen, indem durch Wärmefluss die Temperaturen der beiden Stoffe angeglichen werden, dabei fließt die Wärmemenge \(Q\) immer vom Stoff mit der höheren Temperatur zur niedrigeren. Problem ist nur, dies dauert sehr lange und ist verfälscht durch die Umgebung, es ist kein perfektes thermisch isoliertes System. Stets wird von der Umgebung Wärme abtransportiert.\\
Das heißt, das fl.Wasser kühlt sich ab und erwärmt dadurch das Eis. Es kann jetzt natürlich sein, dass sich eine Mischungstemperatur \(<0\unit{\degreeCelsius}\) einstellt, das würde aber bedeuten, dass \textbf{alles} Wasser im Behälter gefriert. Weder dieser Fall, noch dass alles Eis direkt geschmolzen ist, konnte in der kurzen Zeit beobachtet werden. Gehen wir nun davon aus, da noch Eis vorhanden war, dass das umschließende Wasser drumherum dann ziemlich genau die Schmelztemperatur hat, da es selber nicht gefriert, dann sollte die Messung halbwegs stimmen. Für dieses Wasser hat sich ein Gleichgewicht von pimal Daumen \(T=\qty{0(2)}{\degreeCelsius}\) gebildet, da Umgebungswärme direkt wieder durch das Schmelzen von weiterem Eis kompensiert wurde. Wasser in fester und flüssiger Form haben koexistiert.\\
In \cref{fig:wassermeme.png} sieht man wie Wasser geknechtet wird, liegt aber im Auge des Betrachters (es gibt bestimmt noch viel schlimmere Existenzen von Wasser (je nachdem was als \glqq ideal\grqq{}definiert wird))
\begin{figure}[htbp]            
\centering        
\textbf{\textsf{Born to: \hspace{5cm} forced to:}}\\
\vspace{2mm}
\includegraphics[width=14cm]{wassermeme.png}
\caption[\emph{Meme:} Armes \ce{H2O} \autocite{Wasser}]{Armes \ce{H2O} \autocite{Wasser}}
\label{fig:wassermeme.png}
\end{figure}

\emph{Anmerkung}: Diesen Absatz hab ich ganz am Anfang geschrieben, nachdem ich noch nichts ausgewertet habe. Wie man in \cref{table:allewerte} sieht, der erste Messzeitpunkt beim Wasser-Eis-Gemisch, der mit \qty{0}{\degreeCelsius} angedacht war, zeigt sowohl beim Druckwert \(T_p\) nach der verbesserten Eichkurve 2 UND auch nach \(T_R\), dem Pt100 Herstellerwert, dass die Annahme, das Wasser hätte \qty{0}{\degreeCelsius}, wohl falsch ist. Auch meine anfängliche Idee (oben), dass man ja gerade erst geschmolzenes Wasser misst, scheint falsch zu sein. Vielleicht ist die Durchmischung dazu nicht gut genug gewesen, oder das Eis hat das Wasserhahnwasser noch nicht ausreichend abgekühlt.

\subsection{Eichkurve}
Da die Druckwerte in \cref{table:T(p)} nur sehr ungenau eingetragen werden können, die Skalenteilung ist einfach zu klein, könnte man dort noch einen weiteren Fehler zum Ablesefehler an der Temperaturachse hinzufügen. 
Die Werte konnten nur auf eine Genauigkeit von \(\qty{0.5}{\milli\meter}\equiv\qty{2}{\hecto\pascal}\) eingetragen werden. Eigentlich ist jeder Druckwert fehlerbehaftet mit \(\Delta p=\qty{2}{\hecto\pascal}\equiv\frac{1}{10}\,\unit{\milli\meter}\leftarrow\) das ist der Skalawert des Fehlers. Es ist also unmöglich, diesen Fehler im Diagramm einzuzeichnen. Quadratisch addiert mit dem Einzeichfehler von \(p\) ergibt das:
\[\Delta p_{ges}=\qty{3}{\hecto\pascal}\]
Der daraus resultierende Temperaturfehler ergibt sich durch Analysieren der Steigung der \(T(p)\) Gerade, 
\[\Delta T=\Delta p_{ges}\cdot m=\qty{3}{\hecto\pascal}\cdot\frac{100-(-195.8)}{1242-258}\frac{\unit{\degreeCelsius}}{\unit{\hecto\pascal}}=\qty{1}{\degreeCelsius}\]
Diesen Fehler kann man jetzt mit dem Ablesefehler für \(T\) quadratisch addieren und kommt dadurch auf einen realistischeren Fehler. Zeichnerische Ungenauigkeiten sind aber immer noch nicht berücksichtigt :) Wahrscheinlich ist auch der Einzeichfehler von \qty{0.5}{\milli\meter} zu niedrig geschätzt.

\xfigure{htbp}{scale=0.6}{morpheus2.png}{\emph{Meme:} Ach, Zeichnen per Hand\autocite{Morpheus}}{Das schöne Zeichnen mit Hand, gilt auch für die Fehlergerade von \(R(T)\) in \cref{fig:Widerstand.pdf} \autocite{Morpheus}}

\subsection{Ergebnisse}
Die Abweichung von \(T_0=\qty{-286(2)}{\degreeCelsius}\) der ersten Eichgerade vom absoluten konstanten Wert der Thermodynamik: \(T_0=\qty{-273.15}{\degreeCelsius}\) beträgt \(6.4\sigma\), was natürlich jenseits von gut und böse ist. Der Grund hierfür liegt wahrscheinlich an dem ungenau gemischten Wasser, wie oben diskutiert. 
Mit der zweiten Eichgerade erhalten wir dann einen Wert von \(0.8\sigma\). Man könnte die zweite Eichgerade bestimmt noch genauer machen, indem man den gemessenen Druck der Siedetemperatur von Stickstoff, \(T_{\ce{N2}}=\qty{-195.8}{\degreeCelsius}\) auch noch auf Normalbedingungen umrechnet. Da Stickstoff direkt sublimiert, also die Volumenänderung zwischen beiden Aggregatszuständen sehr groß ist, hat der Umgebungsdruck eine große Auswirkung auf \(T_{\ce{N2}}\).

Die Abweichung des ermittelten Wertes \(T_{\ce{CO2}}=\qty{-62(1)}{\degreeCelsius}\) vom Literaturwert (nicht auf Normalbedingungen umgerechnet, ich weiß nicht wie) \(T_{\text{Lit},\,\ce{CO2}}=\qty{-78.4}{\degreeCelsius}\) beträgt \(16.4\sigma\). Diese hohe Abweichung liegt wahrscheinlich daran, dass bei dieser Versuch der Glasballon nicht richtig in das Dewargefäß eingetaucht werden konnte, dementsprechend hat sich das Gas nicht adäquat abgekühlt, der gemessene Druck ist nämlich viel zu hoch.  

\textbf{Generell} macht die Formel \cref{eq:Druck_Normal} von Jens Wagner, die wir benutzt haben um auf Normalbedingungen umzurechnen keinen Sinn (glaube ich). Denn unser Luftdruck ist niedriger als der Atmosphärendruck. Das heißt, die Stoffe fangen schon bei einer niedrigeren Temperatur an zu sieden. Nach \cref{eq:Druck}, der gemessene Druck eines Gases ist proportional zur Temperatur, müsste unsere angeglichener Druck eigentlich niedriger werden. Durch den Bruch in Wagners-\cref{eq:Druck_Normal} wird er aber größer?!

Würde ich also den gemessenen Druck von \(T_{\ce{CO2}}=\qty{695}{\hecto\pascal}\) nach seiner Formel angleichen, bekommen wir eine noch höhere Sigma-Abweichung (die neue Temperatur ist höher, da \(p_{NB}\) höher wird).
 
\subsection{Widerstandstemperaturen nach Eichtabelle des Pt100}
Weil es mich interessiert hat, habe ich mal in der Herstellertabelle (\cref{fig:eichtabelle.png}) die Temperaturen \(T_R\) rausgesucht, die den gemessenen Widerständen (ich hab mal den Spannungsmessfehler ignoriert) entsprechen, und ihre Fehler berechnet, mittels:
\begin{equation}
    \Delta T=\qty{0.30}{\degreeCelsius}+0.005T
\end{equation}
Die ungefähren Werte (manchmal passt \(R_i\) nicht perfekt auf ein \(T_i\) aus der Tabelle, also manchmal vielleicht um \(\pm\qty{0.5}{\degreeCelsius}\) daneben) verglichen mit den nach \cref{eq:T(R)} berechneten Werten \(T_{R, ger.}\) sind in der folgenden Tabelle. Dort sind zusätzlich noch die Sigma-Abweichungen zu den in der \cref{table:T(p)}(erste Zeile) zu findenden Werten \(T_p\):
\begin{table}[htbp]
    \centering\makegapedcells
    \caption{Temperaturen des Pt100 verglichen mit den Gasthermo-Temperaturen}
\begin{tabular}{|C|*{11}{C}|}
\hline
 R\,[\unit{\ohm}]& 102 & 103,9 & 108,3 & 112,3 & 117 & 120,1 & 124 & 128 & 131 & 135  & 138 \\ \hline
 T_R\,[\unit{\degreeCelsius}]& 5   & 10    & 21    & 32    & 44  & 53    & 62  & 73  & 80  & 91   & 99  \\ \hline
 T_{R, ger.}\,[\unit{\degreeCelsius}]&5	&10	&21	&32	&44	&52&62&	72	&80&	91	&99\\ \hline
 \Delta T_R\,[\unit{\degreeCelsius}]& 0,3 & 0,4   & 0,4   & 0,5   & 0,5 & 0,6   & 0,6 & 0,7 & 0,7 & 0,8  & 0,8 \\ \hline
 T_p\,[\unit{\degreeCelsius}]& 5   & 14    & 26    & 35    & 44  & 55    & 65  & 73  & 82  & 90   & 100 \\ \hline
 \sigma_{abw}& 0   & 2     & 2,5   & 1,5   & 0   & 1     & 1,4 & 0   & 0,9 & 0,5 & 0,5 \\ \hline
\end{tabular}
\label{table:allewerte}
\end{table}

\newpage
Wir sehen zwei sehr schöne Dinge:
\begin{enumerate}
    \item Die berechneten Werte für \(T_R\) sind sehr ähnlich zu den in der Tabelle abgelesenen Werten (deswegen würde beim Auftragen in einem \(R(T)\) Diagramm auch \(R_0A\) als Steigung rauskommen, da \cref{eq:T(R)} auf \cref{eq:R(T)} basiert.) Die Werte der \(T_R\)s unterscheiden sich ziemlich sicher höchstens mal, weil ich in der Tabelle einen evtl. falschen Wert abgelesen habe. Jetzt ist natürlich die Frage, hat der Hersteller die Tabelle auch anhand dieser Formel \cref{eq:R(T)} erstellt, oder experimentell getestet?, bzw. wsl basiert die Formel auch auf experimentellen Werten.
    \item Die Sigma-Abweichungen zu \(T_p\) sind alle akzeptabel, das heißt unsere Eichkurve war gar nicht mal so schlecht. (Vielleicht stimmt der Zeichnen-Meme dann doch nicht?)
\end{enumerate}
Wie bei der Sigma-Abweichung üblich: Würde man die Fehler höher abschätzen, z.B. durch Einbeziehen des Voltmeter-Fehlers (ein Fehler durchs Voltmeter von \(\Delta R\approx \qty{2.7}{\ohm}\) entspricht ungefähr \(\Delta T_V=\qty{7}{\degreeCelsius})\) und mit \(\Delta T_R\) quadratisch addieren, dann hätten wir eine \emph{sehr} kleine Sigma-Abweichung. Wäre bisschen schwachsinnig bzw. nichtsaussagend dann, deswegen habe ich diesen Fehler weggelassen.


\subsection{Temperatur von Stickstoff und CO2 mit dem Pt100}
Auch aus Interesse:
Aus der Herstellertabelle \cref{fig:eichtabelle.png} wurde aus den gemessenen Widerstandswerten die entsprechende Temperatur abgelesen.
\begin{align*}
    R_{\ce{CO2}}=\qty{68.9}{\ohm}&&R_{\ce{N2}}=\qty{19.9}{\ohm}
\end{align*}
Es ergibt sich:
\begin{table}[htbp]
    \centering\makegapedcells
    \caption{Vergleich der tiefen Temperaturen}
\begin{tabular}{|C|CCC|}
\hline
 &  T_{\text{Lit.}}\;[\unit{\degreeCelsius}]    & T_R\,[\unit{\degreeCelsius}]   &   T_p\,[\unit{\degreeCelsius}]   \\ \hline
 \ce{CO2} & -78.4  & -79  & -60  \\ \hline
 \ce{N2} & -195.8 & -199 & -206 \\ \hline
\end{tabular}
\end{table}


Die Abweichungen vom Lit.-Wert sind also für beide Stoffe deutlich kleiner als durch die Druckmessung. Bei \ce{CO2} liegt es an der schwierigen Immersion des Glasballons, bei Stickstoff haben wir wie Gimli \cref{fig:gimli.png} dann genügend \ce{N2} in das Dewargefäß gemacht, verwunderlich ist hier, dass der gemessene Druck auf eine noch niedrigere Temperatur als den Literaturwert zeigt. Ist bei dieser tiefen Temperatur das Gasgesetz nicht mehr näherungsweise gültig? Wahrscheinlich war der Stickstoff aber auch einfach kälter als die Siedetemperatur.


\subsection{Fazit}
All in all eigentlich ganz nett der Versuch, hat aber hauptsächlich in der Auswertung wegen den Memes Spaß gemacht. Das Nette am Versuch ist es, neue Geräte kennenzulernen: Interessant wäre eigentlich vielmehr, die Hintergründe vom piezoresistiven Effekt zu lernen, vom Seebeck-Effekt, dessen Anwendung und mehr zu Wärmestrahlung. Aber gut, die gewollten Messergebnisse konnten ziemlich gut erhoben und ausgewertet werden, in der Hinsicht ist der Versuch erfolgreich verlaufen.
